% Section V -- Verification and validation (WP12 draft, stage C). Every number is a macro of
% numbers.yaml filled by tools/gen_results.py from committed result files (or a code or
% definition macro); figures F3-F6 come from figures/src/*.py. Honesty rules (DESIGN 1.6):
% verification != validation; MATLAB is code-to-code; digitized data are DIG; the lab study
% is a "consistency check under geometric uncertainty" (no lab geometry supplied, H-1).
% The tick-truncation pitfall is a generic design lesson. Tang et al. (V-C) is reported as
% accuracy, measured minus emulator (author decision D-12 of 2026-10-11, reversing DESIGN 10.8).
\section{Verification and Validation}\label{sec:vv}

This section keeps two questions apart. \emph{Verification} asks whether the code computes
what the model states. We answer it with closed forms, exact properties and an independent
code. \emph{Validation} asks whether the model describes the physical world. That needs
measurements. Agreement with MATLAB is therefore verification. Data read off a
published figure are labelled DIG and carry their digitization uncertainty, and every figure
caption names the classes of the data it shows (Appendix~\ref{app:provenance}).

\subsection{Analytic oracles}\label{sec:vv-oracles}
% GENERATED by paper/tools/gen_results.py from the test record in generated/junit/
% (sidecar meta: sremu commit 5be649a3a0018c848ffff087a3dfa934997b4f14). Do not edit.
\begin{table*}[t]
\caption{Verification summary from the test record of the emulator commit used for this paper. An identifier passes when at least one implementation test that carries it passed and none failed or was skipped; rendering-path identifiers count once per backend. Tolerances and oracles of every identifier are fixed in the specification; Appendix~\ref{app:oracles} describes the oracle families. PH-12, a regression value of the experimental P.526 module that lies outside every gate, is not counted; its implementation tests run in the experimental test tier, outside the continuous-integration and nightly runs that this record covers.}
\label{tab:vv}
\centering
\footnotesize
\setlength{\tabcolsep}{3pt}
\renewcommand{\arraystretch}{1.12}
\begin{tabular}{@{}>{\raggedright\arraybackslash}p{0.35in}>{\raggedright\arraybackslash}p{3.65in}>{\raggedright\arraybackslash}p{1.95in}>{\raggedleft\arraybackslash}p{0.45in}>{\raggedleft\arraybackslash}p{0.4in}@{}}
\toprule
Tier & What is compared & Identifiers passed & Passed & Tests \\
\midrule
PH & Physics closed forms and exact properties: Friis, area-tied kernel equals published forms, aperture refinement, image theory, focusing, ground, knife edge, alphabets and optimum, Doppler, interpolation, fractional delay, noise, linearity & PH-01--PH-11, PH-13--PH-32 & 31/31 & 352 \\
RT & Runtime: exact counts at several rates, round trips, bitwise invariance, switch samples, late commands, ties, barriers, oscillator continuity, transitions, replay, paced counters, backends & RT-01--RT-19 & 19/19 & 192 \\
IO & Interfaces: malformed frames and messages, two-process bitwise equality, socket commands, caps and loopback, SigMF, status codes, wrappers, raw mode & IO-01--IO-08 & 8/8 & 127 \\
ST & Statistical: random-waypoint law, diffuse autocorrelation and level crossings, quantization and failure Monte Carlo, BER vs theory, CFO vs CRB & ST-01--ST-05 & 5/5 & 30 \\
G & MATLAB R2026a code-to-code: per-element kernel, beam patterns, CFO, thermal noise, fading statistics, ETSI far field, two-ray channel & G1--G8 & 8/8 & 30 \\
\bottomrule
\end{tabular}
\end{table*}

\claim{C-VV-TABLE}Table~\ref{tab:vv} summarizes the test record of the emulator commit used
for this paper. Each identifier names one property, its oracle and a fixed tolerance. The
oracles are closed forms or exact properties. They live in a separate test package that may
not import the code under test (Appendix~\ref{app:oracles}). All \numOracleCount{} analytic
identifiers (PH, RT and IO) pass, in \numOracleTests{} implementation tests. The statistical
tier adds \numStatIds{} identifiers and the code-to-code tier \numGoldenIds{}. In total
\numJunitPassed{} tests passed and \numJunitSkipped{} were skipped. No skipped test carries an
identifier: they test fallback behaviour for modules that are now present, and a
processor-affinity call that this platform lacks. An identifier counts as passed when at
least one of its implementation tests passed and none failed or was skipped. Tests of the
rendering path count once per backend. Appendix~\ref{app:provenance} lists the commit and
date of this record with those of every other result; the record does not state its host.

The analytic tier spans the whole emulator. On the physics side it checks the Friis
equation, the equality of K1 with three published forms, image theory, the $N^2$ law of a
focused panel, the two-ray model, the knife edge and the losses of finite alphabets. On the
runtime side it checks exact sample counts at incommensurate rates, the switch sample of a
distant receiver, late and tied commands, barrier determinism and bitwise invariance under
any chunking of the client calls. Section~\ref{sec:design} cites these identifiers where it
states a property.

\subsection{MATLAB code-to-code comparison}\label{sec:vv-matlab}
\claim{C-MATLAB}Eight golden references, G1--G8, come from MATLAB R2026a scripts. G1 sums K1
per element with toolbox primitives for laboratory-like layouts and for the geometries of
Tang et al. G2 and G3 compare steered and one-bit beam patterns with MATLAB's RIS reflector
object. G4, G5 and G6 compare the frequency offset, thermal noise and fading statistics with
communications-toolbox objects. G7 evaluates the far-field closed form of ETSI GR RIS
003~\cite{ETSI_GR_RIS003_v121}, and G8 is MATLAB's two-ray channel. All \numGoldenIds{} pass
(Table~\ref{tab:vv}). MATLAB is not needed to run them: the outputs are stored with the
script hash and the MATLAB version. These comparisons show that two codes agree; they are not
validation. MATLAB's reflector uses unit-gain cells, so G2 and G3 compare beam shapes after a
documented constant is removed. They say nothing about absolute RIS power.

\subsection{Published RIS path-loss model and measurements}\label{sec:vv-tang}
\claim{C-REFINE}Fig.~\ref{fig:F3}(a) checks a property that published curves cannot test.
When a fixed aperture is cut into more cells, the received power must not change. Far from
the panel, at specular reflection, K1 with the default cell model stays within
\numRefineFarMaxDevDb~dB of the closed form for every cell count from \numRefineNmin{} to
\numRefineNmax{} (PH-04). Near the panel, at an anomalous angle, the patch model stays within
\numRefineNearMaxDevDb~dB of the continuous physical-optics integral, which we evaluate
independently. One point per cell is off by \numRefinePointCoarseDb~dB for the coarsest
cells and converges slowly. At the smallest count with a pitch below half a wavelength
(\numRefinePointFirstAutoN{} cells) it is still off by \numRefinePointFirstAutoDb~dB. This is
the model-form difference of Section~\ref{sec:model-cells}, not a refinement error, and it is
why the default switches to patches for larger cells.

\claim{C-ABS-TANG}Tang et al.\ measured the path loss of millimeter-wave surfaces in an
anechoic chamber and fitted a model to it~\cite{Tang2022mmWave}. We extracted the model curves
and the measured points of their Fig.~10 from the vector graphics of the
paper~\cite{Tang2022mmWave} (DIG) and ran the same \numTangMeasN{} links, four geometries
of nine receiver distances, through the emulator with the default cell model, the reported
unit-cell state and the reported link budget (REPORTED). The emulator reproduces the
published model curves within \numTangCurveMaxDevDb~dB; since K1 equals that model for one
point per cell (Section~\ref{sec:model-k1}), this checks the implementation. Against the
chamber measurements, the emulator's received power is low by \numTangMeasBiasDb~dB on
average (SD \numTangMeasSdDb~dB, RMSE \numTangMeasRmseDb~dB, largest difference
\numTangMeasMaxDb~dB; per geometry \numTangPanelBiasMinDb--\numTangPanelBiasMaxDb~dB,
Appendix~\ref{app:tang}), and every measured point lies above it (Fig.~\ref{fig:F3}(b)). This is
the accuracy of the model class that the emulator shares with Tang et al.'s model, against
third-party measurements of one \numTangFcGHz~GHz surface in one state and in specular
geometries: it tests the absolute level and its dependence on distance and angle, not
reconfiguration or near-field structure. The points are digitized, the source states no
measurement uncertainty, their distances were snapped to the nominal grid of receiver
distances, and the points of one geometry are not independent. We do not fit
away the constant part of the error, which may come from the reported link-budget terms.

\begin{figure*}[t]
\centering
\sremufigure{F3}{\textwidth}{(a) Fixed-aperture refinement of K1: deviation of the received
power from a reference as one aperture is cut into $N$ cells. Top: far specular reflection,
default cell model, against the far-field closed form. Bottom: near-field anomalous
reflection, patch model and a forced point model, against the continuous physical-optics
integral. (b) Published model curves and digitized chamber measurements of Tang et
al.~\cite{Tang2022mmWave} and the emulator's values for the same links, four geometries.
Data classes: ANALYTIC (references in (a)), SIM-sremu (emulator), DIG (published curves and
measurements), REPORTED (link budget).}{fig:F3}
\end{figure*}

\subsection{Laboratory consistency check}\label{sec:vv-lab}
\claim{C-LAB-DATA}The data set is a laboratory measurement collected by the
authors: the received power at \numLabFreqGHz~GHz on a polar grid of \numLabCells{} cells in
front of a binary RIS of six tiles. What survives is a colour map, so we digitized it
(DIG, \numLabSigDigMinDb--\numLabSigDigMaxDb~dB per cell). The geometry was not recorded to
the precision a model needs. The distance from the transmitter to the tiles, the angular span
of the tiles, the cell type and the measurement positions are known within ranges, the
positions to about $\pm\numLabPriorPosCm$~cm.

\claim{C-LAB-PRESPEC}We therefore wrote the analysis plan before any comparison of model and
data. The plan draws \numLabDraws{} geometries from a prior over the unknown quantities and
predicts the map with five models, each with a fitted offset: line of sight alone (M0),
a log-distance fit that also fits its exponent (M0n), a uniform mirror (M1), the laboratory configuration (M2) and M2 with
a floor (M3). The offsets are fitted on three rings (\numLabTrainCells{} cells) and scored on
the other two (\numLabTestCells{} cells). Three decision rules were fixed in advance. C1
(level consistent) needs a median absolute bias of at most \numLabConeBiasMaxDb~dB and a
median RMSE of M2 at most \numLabConeMarginDb~dB above that of M0n. Because each offset absorbs the unknown power and
gains, C1 tests the level of the held-out rings, not absolute power. C2 (the RIS model
explains spatial structure) needs a median RMSE gain of M2 over M0 of at least
\numLabCtwoMinDb~dB, a paired-bootstrap interval that excludes zero in at least
\numLabCtwoCiFracPct~\% of the draws, and a Spearman correlation of at least
\numLabCtwoRhoMin{} whose interval excludes zero. C3 is C2 with M1 in place of M0. The protocol
files were hashed, and the data loader refuses the measured map unless the hash matches. The
hash was committed to the emulator repository (commit \txtLabLockCommit) and pushed to its
private GitHub repository on \txtLabLockDate, after a descriptive inspection of the
digitized map but before any comparison of model and data. GitHub's server-side push log,
exported later, records that push (Appendix~\ref{app:provenance}). Two commits followed: an
option of the run script that names the columns of the measured file, and a test change
(commit \txtLabCodeCommit), pushed at \txtLabCodePushTime; the unblinded run started at
\txtLabRunTime{} from a clean tree of that commit, and its result
(commit \txtLabResultCommit) was pushed at \txtLabResultPushTime. No public registry deposit
(Zenodo or OSF) was made before unblinding; the plan, its hash and the lock records can still
be deposited, as a retrospective record. The plan is therefore pre-specified and hash-locked,
with a timestamp from the repository host, not from a public registry.

\claim{C-LAB-POWER}The plan also asked whether the test could succeed at all. Before
unblinding, it ran the full analysis on \numLabDecoys{} synthetic maps drawn from M2 itself.
C2 held in \numLabDecoyDetections{} of them (\numCiPct\,\% upper limit \numLabDecoyCiHighPct~\%),
a detection rate far below the pre-set threshold of \numLabUnderpoweredRate{}. C1 held in
\numLabDecoyConeDetections{} and C3 in \numLabDecoyCthreeDetections{} of them. The study was labelled
underpowered before the data were opened.

\claim{C-LAB-RESULT}The outcome is a null result (Table~\ref{tab:lab}, Fig.~\ref{fig:F4}). The
median test RMSE is \numLabRmseMzero~dB for M0, \numLabRmseMzeroN~dB for M0n and
\numLabRmseMtwo~dB for M2 (\numQuantLoPct--\numQuantHiPct\,\% range:
\numLabRmseMtwoLo--\numLabRmseMtwoHi~dB), with a
median M2 bias of \numLabBiasMtwo~dB. The RMSE gain of M2 over M0 is \numLabDeltaRmse~dB, its
bootstrap interval excludes zero in \numLabCiExcludesPct~\% of the draws, and the Spearman
correlation is \numLabSpearman{} (\numCiPct\,\% interval \numLabSpearmanLo{} to
\numLabSpearmanHi). C1, C2 and C3 all fail, also in every pre-set sensitivity run. For
orientation, outside the decision rules: the observed RMSE gain lies close to the median
of the decoys drawn from M2 (\numLabDecoyDeltaMtwoDb~dB) and above that of decoys drawn from
M0 (\numLabDecoyDeltaMzeroDb~dB). The digitized map does not discriminate the RIS model from LOS-only at the available
geometric knowledge. The comparison is underpowered, so this outcome is not evidence for or
against the model. One cell, pinned to the edge of the colour scale, has a residual of
\numLabResidualOutlierDb~dB.

\claim{C-LAB-LESSON}Which extra fact would have made the test informative? A diagnostic
outside the protocol repeated the power analysis with narrower priors. None of
\numLabDiagSingleFacts{} single facts (for example the transmitter-to-tile distance to a few
millimetres, the tile angles, the cell type, the antenna pattern or the state reflection
coefficients, or even the measurement positions alone to $\pm\numLabDiagPosMm$~mm) lifted
the C2 detection rate above zero. Knowing the as-built tile geometry
exactly and the measurement positions to $\pm\numLabDiagPosMm$~mm gave
\numLabDiagRate{} (\numCiPct\,\% interval \numLabDiagRateLo--\numLabDiagRateHi, \numLabDiagDecoys{} decoys);
with positions to $\pm\numLabDiagPosFiveMm$~mm it stayed at zero, and knowing every
geometric quantity exactly gave \numLabDiagGeomRate. At \numLabFreqGHz~GHz the wavelength is
\numLabLambdaMm~mm, and the map's structure comes from path differences of a fraction of it.
A RIS map at this frequency therefore needs its geometry recorded to about a millimetre, for
example with a fixture or a laser tracker, before it can separate the models. Even then the
detection rate stays close to the pre-set threshold. We report this as a lesson for
measurement design; the word \emph{validation} is reserved for a data set with a recorded
geometry.

% GENERATED by paper/tools/gen_results.py from sremu validation/ants_fig4b/results/20261011T003044Z_vlab_unblinded/table_v.json.
% Do not edit.
\begin{table}[t]
\caption{Laboratory consistency check: test-ring metrics over the prior draws of the pre-registered protocol (median, with 5--95\,\% range for the RMSE). Error = measured $-$ predicted. Measured: DIG; predictions: SIM-sremu; geometry and $|\Gamma|$: ASSUMED (prior).}
\label{tab:lab}
\centering
\scriptsize
\setlength{\tabcolsep}{2.5pt}
\begin{tabular}{@{}lrrrr@{}}
\toprule
Model & RMSE (dB) & Bias (dB) & MAE (dB) & max$|e|$ (dB) \\
\midrule
M0 (line of sight) & 3.96 (3.92--3.99) & \ensuremath{-0.59} & 2.18 & 12.6 \\
M0n (log-distance) & 3.77 (3.73--3.81) & \ensuremath{-0.47} & 2.22 & 11.9 \\
M1 (mirror) & 5.00 (3.36--6.75) & \ensuremath{-0.64} & 3.50 & 13.5 \\
M2 (lab configuration) & 5.11 (3.27--6.66) & \ensuremath{-0.60} & 3.53 & 13.9 \\
M3 (M2 + floor) & 5.18 (3.30--6.79) & \ensuremath{-0.63} & 3.60 & 14.1 \\
\midrule
\multicolumn{5}{@{}p{3.3in}@{}}{Decisions: C1 fails; C2 fails; C3 fails; decoy C2 detection rate 0.00 (underpowered).} \\
\bottomrule
\end{tabular}
\end{table}

\begin{figure*}[t]
\centering
\sremufigure{F4}{\textwidth}{Laboratory consistency check under geometric uncertainty.
(a) Digitized map of a laboratory measurement collected by the authors. (b) Median
prediction of the laboratory-configuration model M2 over the prior draws. (c) Measured minus
predicted; the colour scale is clipped and the outlier is marked. Test-ring cells are
outlined. (d) RMSE gain of M2 over the line-of-sight model M0 on the test rings over the prior
draws, with the pre-specified C2 threshold. Data classes: DIG (measured), SIM-sremu
(predictions), ASSUMED (geometry and reflection-coefficient prior, decision rules).}{fig:F4}
\end{figure*}

\subsection{CFO estimation and link metrics}\label{sec:vv-cfo}
\claim{C-CFO}A receiver oscillator offset $f$ rotates the received baseband at $f$. With noise
off we swept $f$ over \numCfoSweepPoints{} values within $\pm\numCfoSweepMaxHz$~Hz on a BPSK
burst of \numCfoBurstSamples{} samples at \numCfoRateSps~S/s. A least-squares phase estimator
recovers $f$ within \numCfoLsMaxErrHz~Hz and a lag-1 estimator within \numCfoLagMaxErrHz~Hz
(Fig.~\ref{fig:F5}(a)). The received samples are identical for three different client epoch
lengths (V-CFO-1). The same sweep through two
operating-system processes over the native ZMQ data plane gives bitwise identical samples
(V-CFO-3).

\claim{C-CFO-LESSON}Fig.~\ref{fig:F5}(a) also shows a generic pitfall. At
\numCfoRateSps~S/s a tick of \numCfoTickMs~ms holds \numCfoSamplesPerTick{} samples. A time
base that delivers a whole number of samples per tick, here \numCfoSamplesDelivered{}, drops
samples without notice, and every frequency estimate is scaled by the ratio of the two counts:
a slope of \numCfoTickScale{}, an error of \numCfoTickScalePct~\%. We reproduced the effect by
re-timing the emulator's own output. Exact rational sample counts
(Section~\ref{sec:design-time}) remove it.

\claim{C-CRB}With the emulator's thermal noise at per-sample SNRs of \numCfoSnrLo{},
\numCfoSnrMid{} and \numCfoSnrHi~dB and \numCfoTrials{} keyed trials at each of
\numCfoPoints{} offsets, a periodogram estimator on \numCfoHeaderSamples{} samples reaches the
Cram\'er--Rao bound of a tone in white Gaussian noise~\cite{Kay1989Freq}. Its pooled RMSE is
\numCfoRatioZero, \numCfoRatioTen{} and \numCfoRatioTwenty{} times the bound. An unbiased
estimator cannot beat the bound, so values below one are sampling error of the pooled
estimate. The lowest lies \numCfoRatioTwentySe{} standard errors below one. We checked that the
emulator's noise is white and has the stated power, and the gate does not depend on this
value. The worst single offset at the two gated SNRs is \numCfoRatioMaxPoint{} times the bound, within the
pre-set limit of \numCfoGateFactor{} (Fig.~\ref{fig:F5}(b); V-CFO-4, ST-05). The lag-1
estimator follows its own small-error theory but not the bound: its variance falls with the
number of samples, the bound's with its cube.

\claim{C-BER}Fig.~\ref{fig:F5}(c) sends BPSK and QPSK through a static lab-like link with the
emulator's physical noise. The transmit power sets $E_b/N_0$ from \numBerEbnoMin{} to \numBerEbnoMax~dB, and
each point collects at least \numBerMinErrors{} errors or \numBerMaxBits{} bits. With genie
channel knowledge the theoretical bit error rate lies inside the \numCiPct\,\% Clopper--Pearson
interval at all \numBerPoints{} points. With a least-squares channel estimate from
\numBerPilots{} pilots, theory falls outside the interval at \numBerLsOutside{} points, all at
$E_b/N_0\le\numBerLsOutsideMaxEbnoDb$~dB, as expected from the estimation loss. ST-04 repeats
the genie check with \numBerStCiPct\,\% intervals in the test record.

\begin{figure*}[t]
\centering
\sremufigure{F5}{\textwidth}{(a) Error of the estimated frequency offset versus the applied
offset, noise off: the emulator (exact sample counts) and a generic pitfall, the emulator's own
output re-timed with a whole number of samples per tick, with its analytic slope. (b) RMSE of
two estimators versus SNR against the Cram\'er--Rao bound and the lag-1 small-error theory.
(c) BER versus $E_b/N_0$ through the emulator with genie and pilot-based detection,
Clopper--Pearson intervals, and theory; markers are offset horizontally for legibility.
Data classes: SIM-sremu (emulator runs), ANALYTIC
(bound, theory, slope), ASSUMED (unit amplitudes of the lab states in (c)).}{fig:F5}
\end{figure*}

\subsection{Doppler and mobility}\label{sec:vv-doppler}
\claim{C-DOPPLER-VV}Doppler shifts follow from the geometry (Section~\ref{sec:model-doppler}).
In Fig.~\ref{fig:F6}(a) a receiver approaches the transmitter at \numDopVms~m/s at
\numDopFcGHz~GHz and passes a panel. The rendered line-of-sight path has the closed-form
shift of \numDopLosHz~Hz within \numDopLosDevHz~Hz over one second (PH-17). The RIS
contribution, taken as the difference between runs with and without the panel, follows an
independent per-cell sum of K1. Near a null of the panel pattern the phase of the sum jumps
and no frequency is defined. We therefore compare the samples where the reference RIS
amplitude exceeds \numDopRisNullPct~\% of its maximum (\numDopRisKeptPct~\% of the samples).
There the largest deviation is \numDopRisDevHz~Hz and the median \numDopRisMedHz~Hz. The interpolation bound
applies to the whole channel, so a weak RIS term inherits a larger relative error.

Fig.~\ref{fig:F6}(b) shows the error of the snapshot grid. At the automatic period of
\numDopTuAutoUs~$\mu$s the largest relative error against a channel evaluated at every sample
is \numDopErrAuto, well below the target $\varepsilon_{\mathrm{rel}}=\numEpsRel$. The error
grows as $T_u^2$ and reaches \numDopErrMax{} at \numDopTuMaxMs~ms. The engine's own
estimate, \texttt{verify\_interp}, follows it but reads low, by up to \numDopViMaxLowPct~\% at
the largest period. PH-18 checks the bound against per-element brute
force in three scenes on both backends.

\begin{figure}[t]
\centering
\sremufigure{F6}{\columnwidth}{(a) Doppler shift of the rendered line-of-sight path and RIS
contribution of a moving receiver against the closed forms; RIS samples near nulls of the
panel pattern are omitted. (b) Largest relative interpolation error over one second versus
the snapshot period: against a channel evaluated at every sample, and as estimated by the
engine. Data classes: engine runs and closed forms with the target (SIM-sremu;
ANALYTIC).}{fig:F6}
\end{figure}
